\documentclass[11pt, a4paper]{article}
\usepackage[utf8]{inputenc}
\usepackage{amsmath, amssymb}
\usepackage{geometry}
\geometry{margin=1in}
\usepackage{hyperref}
\usepackage{booktabs}

\title{\textbf{Read-Path Refusal Specification (RPRS-1)}\\ Version 2.0.0}
\author{Sciencedelic Metatech}
\date{\today}

\begin{document}

\maketitle

\begin{abstract}
The Read-Path Refusal Specification (RPRS-1) defines the normative contract for a conforming memory system. A conforming system MUST return either the exact stored value or an explicit error. It MUST NOT return a value whose integrity verification has failed. RPRS-1 v2.0.0 extends the baseline specification to include \textbf{Portable Proofs}, \textbf{Cross-Anchor Monotonicity}, and \textbf{Erasure Resurrection (T10)} protection, eliminating vendor lock-in while preserving cryptographic truth.
\end{abstract}

\section{Scope and Conformance}
A system is RPRS-1 v2.0.0 conformant if it implements the core contract (Section 2), the four cumulative levels (Section 3), the ten attack classes (Section 4), and the portable proof extension (Section 5).

\section{The Core Contract}
For any key \(k\), the read function MUST return:
\begin{equation}
    \texttt{get}(k) = v \quad \lor \quad \texttt{Err(Tamper)} \quad \lor \quad \texttt{Err(Missing)}
\end{equation}
There is no third option. A silent-wrong answer is a critical failure.

\section{Conformance Levels}
The levels are cumulative.
\begin{itemize}
    \item \textbf{Level L1 (Baseline):} Cryptographic tags on all records. Detects payload tamper.
    \item \textbf{Level L2 (Audit):} Appends to a transparency log. Detects container rollback post-facto.
    \item \textbf{Level L3 (Read-Path):} Verifies tags and log before serving. Refuses to serve tampered or rolled-back data.
    \item \textbf{Level L4 (Anchored):} Integrates an external high-water-mark anchor (e.g., S3 Object Lock) to prevent log+container rollback.
\end{itemize}

\section{Attack Classes (T1--T10)}
A conforming L4 system MUST reject all ten storage-level attacks on the read path.
\begin{table}[h]
\centering
\begin{tabular}{@{}lll@{}}
\toprule
\textbf{ID} & \textbf{Attack Class} & \textbf{Required Defeat Mechanism} \\
\midrule
T1 & Content tamper & Record HMAC tag \\
T2 & Tail truncation & Log chain / Append-only \\
T3 & Middle deletion & Log chain \\
T4 & Reordering & Monotonic version check \\
T5 & Forged insertion & Record HMAC tag \\
T6 & Cross-context replay & Key binding / Substitution resistance \\
T7 & Rollback replay & \texttt{max\_version} + \texttt{per\_key\_max} \\
T8 & Metadata tamper & HMAC covers all headers \\
T9 & Snapshot rollback & External anchor (L4) \\
\textbf{T10} & \textbf{Erasure resurrection} & \textbf{Tombstones + Portable Proofs} \\
\bottomrule
\end{tabular}
\end{table}

\subsection{T10: Erasure Resurrection}
An attacker attempts to restore a record that was legitimately deleted by importing an older state. To defeat T10, a conforming system MUST maintain signed tombstones. When importing or reading state, if a tombstone exists for key \(k\) with version \(v_{tomb}\), any record for \(k\) with version \(v_{rec} < v_{tomb}\) MUST be rejected.

\section{Portable Proof Extension}
A conforming system MUST NOT rely on vendor lock-in for its security. It MUST allow the export and import of the entire verified state.

\subsection{Bundle Format}
A Portable Bundle is a single file containing:
\begin{itemize}
    \item The Genesis state (or a reference to it).
    \item All active records and tombstones.
    \item The transparency log tip.
    \item An anchor receipt (optional but required for L4).
    \item A Bundle MAC signed by the container key.
\end{itemize}

\subsection{Import Verification}
Before importing a bundle, the system MUST verify:
\begin{enumerate}
    \item The Bundle MAC.
    \item The internal HMAC tags of every record.
    \item Log monotonicity (no duplicate versions, no gaps).
    \item Cross-anchor monotonicity (the bundle's version MUST be \(\ge\) the current anchor's version).
\end{enumerate}
If any check fails, the system MUST raise \texttt{TamperError} or \texttt{RollbackError} and refuse the import.

\subsection{Cross-Anchor Ceremony (Reanchor)}
To prevent lock-in, a system MUST allow migration between anchors. The \texttt{reanchor(new\_anchor)} function MUST:
\begin{enumerate}
    \item Verify the current log tip matches the old anchor's latest receipt.
    \item Have the new anchor sign the old tip hash.
    \item Begin the next epoch bound to the new anchor.
\end{enumerate}
This allows an enterprise to move from AWS to GCP without breaking the chain of trust.

\section{Status Taxonomy}
A conforming system MUST return one of the following statuses:
\begin{itemize}
    \item \texttt{ok}: Value returned, MAC verified.
    \item \texttt{missing}: Key never existed.
    \item \texttt{tampered}: MAC failed.
    \item \texttt{revoked}: Record intentionally revoked by an admin.
    \item \texttt{tombstoned}: Record intentionally deleted.
    \item \texttt{expired}: Record past its TTL.
\end{itemize}

\section{Security Considerations}
The system MUST assume an attacker with write access to the storage medium but without the encryption keys. The keys MUST reside outside the data file (out-of-band).

\end{document}